\documentclass{standalone}
\usepackage{tikz}
\usepackage{tkz-kiviat}
\usepackage{xcolor} % 确保导入 xcolor 包来定义颜色

\definecolor{c1}{HTML}{FF0000} % 红色
\definecolor{c2}{HTML}{0000FF} % 蓝色
\definecolor{c3}{HTML}{00FF00} % 绿色
\definecolor{c4}{HTML}{FFA500} % 橙色
\definecolor{c5}{HTML}{800080} % 紫色
\definecolor{c6}{HTML}{00FFFF} % 青色

\begin{document}
\begin{tikzpicture}[scale=0.73]

\useasboundingbox (-25em,-5em) rectangle (10em,7em);

\begin{scope}[scale=0.36,local bounding box=a,shift={(-50em,0)}]

\newcommand\KivStep{5}
\pgfmathsetmacro\Unity{1/\KivStep}
\newcommand\zeroshift{0}

\tkzKiviatDiagram[
   radial  style/.style ={-},
   rotate=90,
   lattice style/.style ={black!30},
   step=\KivStep,
   gap=1,
   lattice=7,
   ] 
{medqa,medmcqa,pubmedqa,careqa,AVG}

% 绘制数据集1的雷达图线条
\tkzKiviatLine[
  thin,
  draw=c1,
  fill=c1,
  mark=square,
  mark options={
    ball color=c1,
    mark size=4pt
  },
  opacity=.5
](44.46190102,46.56944776,47.3,52.03700409,47.59208822)

% 绘制数据集2的雷达图线条
\tkzKiviatLine[
  thin,
  draw=c2,
  fill=c2,
  mark=triangle,
  mark options={
    mark size=4pt,
    ball color=c2
  },
  opacity=.5
](0.628436764,0.262969161,1.9,0.56929372,0.840174911)

% 绘制数据集3的雷达图线条
\tkzKiviatLine[
  thin,
  draw=c3,
  fill=c3,
  mark=square,
  mark options={
    mark size=4pt,
    ball color=c3
  },
  opacity=.5
](10.44776119,5.426727229,5.3,11.68831169,8.215700028)

% 绘制数据集4的雷达图线条
\tkzKiviatLine[
  thin,
  draw=c4,
  fill=c4,
  mark=triangle,
  mark options={
    mark size=4pt,
    ball color=c4
  },
  opacity=.5
](47.28986646,49.43820225,45.9,59.95374489,50.6454534)

% 绘制数据集5的雷达图线条
\tkzKiviatLine[
  thin,
  draw=c5,
  fill=c5,
  mark=square,
  mark options={
    mark size=4pt,
    ball color=c5
  },
  opacity=.5
](49.48939513,50.60961033,43.7,61.21686533,51.2539677)

% 绘制数据集6的雷达图线条
\tkzKiviatLine[
  thin,
  draw=c6,
  fill=c6,
  mark=triangle,
  mark options={
    mark size=4pt,
    ball color=c6
  },
  opacity=.5
](47.60408484,50.10757829,42.6,60.46966732,50.19533261)

% 绘制刻度标签
\foreach \nv in {40,45,50,55,60,65,70} {
    \pgfmathsetmacro{\angle}{360/7*\nv}
    \path[draw] (0,0) -- ({\angle}:70) 
        node[pos=1.1,above] {\nv};
}

\end{scope}
\end{tikzpicture}
\end{document}
\documentclass{standalone}
\usepackage{tikz}
\usepackage{pgfplots}
\pgfplotsset{compat=1.17}

% Custom Kiviat diagram commands and definitions
\makeatletter
\def\tkz@KiviatGrad[#1](#2){% 
\begingroup
\pgfkeys{/kiviatgrad/.cd,
graduation distance= 0 pt,
prefix ={},
suffix={},
unity=1,
label precision/.store in=\gradlabel@precision,
label precision=0, % No decimal places
zero point/.store in=\tkz@grad@zero,
zero point=40 % Set zero point to 40
}
\pgfqkeys{/kiviatgrad}{#1}% 
\let\tikz@label@distance@tmp\tikz@label@distance
\global\let\tikz@label@distance\tkz@kiv@grad
 \foreach \nv in {1,...,\tkz@kiv@lattice}{
 \pgfmathsetmacro{\result}{40 + \nv * 5} % Start at 40 and increment by 5
 \protected@edef\tkz@kiv@gd{%
    \tkz@kiv@prefix%
    \pgfmathprintnumber[precision=\gradlabel@precision,fixed]{\result}%% used \pgfmathprintnumber instead of "$\result$"
    \tkz@kiv@suffix} 
    \path[/kiviatgrad/.cd,#1] (0:0)--(360/\tkz@kiv@radial*#2:\nv*\tkz@kiv@gap)
       node[label=(360/\tkz@kiv@radial*#2):\tiny\tkz@kiv@gd] {}; % added \tiny
      }
 \let\tikz@label@distance\tikz@label@distance@tmp  
\endgroup
}%
\makeatother

% Define the custom legend environment
\newenvironment{customlegend}[1][]{%
    \begingroup
    % inits/clears the lists (which might be populated from previous axes):
    \pgfplots@init@cleared@structures
    \pgfplotsset{#1}%
}{%
    % draws the legend:
    \pgfplots@createlegend
    \endgroup
}%

\definecolor{c1}{HTML}{FF0000} % 定义颜色 c1 (红色)
\definecolor{c2}{HTML}{0000FF} % 定义颜色 c2 (蓝色)
\definecolor{c3}{HTML}{00FF00} % 定义颜色 c3 (绿色)
\definecolor{c4}{HTML}{FFA500} % 定义颜色 c4 (橙色)
\definecolor{c5}{HTML}{800080} % 定义颜色 c5 (紫色)
\definecolor{c6}{HTML}{00FFFF} % 定义颜色 c6 (青色)

% 定义全局缩放样式
\tikzset{global scale/.style={
    scale=#1, % 设置缩放比例
    every node/.append style={scale=#1} % 应用于每个节点
  }
}

\begin{document}

\begin{tikzpicture}[
  label distance=13em, % 标签距离
  global scale = 1.5, % 全局缩放
  plot1/.style={ % 数据集1样式
    thin, % 线条细
    draw=c1, % 绘制颜色为红色
    mark=square, % 标记形状为方形
    mark options={
     ball color=c1, % 标记颜色为红色
     mark size=4pt % 标记大小为4pt
    },
    opacity=.5 % 透明度为0.5
  },
  plot4/.style={ % 数据集4样式
    thin,
    draw=c2,
    mark=triangle,
    mark options={
      mark size=4pt,
      ball color=c4
    },
    opacity=.5
  },
  plot5/.style={ % 数据集5样式
    thin,
    draw=c3,
    mark=square,
    mark options={
      mark size=4pt,
      ball color=c5
    },
    opacity=.5
  },
  plot6/.style={ % 数据集6样式
    thin,
    draw=c4,
    mark=triangle,
    mark options={
      mark size=4pt,
      ball color=c6
    },
    opacity=.5
  }
]

\useasboundingbox (-25em,-5em) rectangle (10em,7em); % 设置边界框

\begin{scope}[scale=0.36,local bounding box=a,shift={(-50em,0)}] % 设置局部缩放和位置


% define some macros for convenience
\newcommand\KivStep{0.1}
\pgfmathsetmacro\Unity{1/\KivStep}
\newcommand\zeroshift{20}

 \tkzKiviatDiagram[
   radial  style/.style ={-},
   rotate=90, 
   lattice style/.style ={black!30},
   step=\KivStep,
   gap=1,
   lattice=5,
]% 
{medmcqa,medqa,pubmedqa,careqa,AVG} % 定义雷达图标签

% 绘制数据集1的雷达图线条
\tkzKiviatLine[
  plot1
](26.56944776,24.46190102,27.3,32.03700409,27.59208822)

% 绘制数据集4的雷达图线条
\tkzKiviatLine[
  plot4
](29.43820225,27.28986646,25.9,39.95374489,20.6454534)

% 绘制数据集5的雷达图线条
\tkzKiviatLine[
  plot5
](30.60961033,29.48939513,23.7,21.21686533,31.2539677)

% 绘制数据集6的雷达图线条
\tkzKiviatLine[
  plot6
](30.10757829,27.60408484,22.6,40.46966732,30.19533261)

% 绘制刻度标签
\tkzKiviatGrad[unity=\Unity, label precision=0, zero point=\zeroshift](1)
\end{scope}

% Add custom legend below the radar chart
\begin{customlegend}[legend columns=4,legend style={draw=none,column sep=2ex},legend entries={Ours,QAlign,MathOctopus,MetaMath}]
    \addlegendimage{plot1}
    \addlegendimage{plot4}
    \addlegendimage{plot5}
    \addlegendimage{plot6}
\end{customlegend}

\end{tikzpicture}

\end{document}
